% Elemento opcional. Insira aqui a errata do documento, caso necessário.
% A errata lista as correções que devem ser feitas no texto após a impressão.
%
% Exemplo de formato:
%
% SOBRENOME, N. \textbf{Título do Trabalho}. Ano. Tipo -- Instituição, Local, Ano.
%
% Deixe este arquivo sem conteúdo (apenas comentários) para omitir a errata.

% [Insira a referência do documento aqui]
SILVA, M. E. \textbf{Análise de um Problema de Pesquisa}: um estudo de caso ilustrativo. 2026. Relatório Técnico --- Embrapa Agricultura Digital, Campinas, 2026.

\begin{table}[htb]
	\centering
	\footnotesize
	\begin{tabular}{|p{1.4cm}|p{1cm}|p{3cm}|p{3cm}|}
		\hline
		\textbf{Folha} & \textbf{Linha}  & \textbf{Onde se lê}  & \textbf{Leia-se}  \\
		\hline
		% [Adicione as correções aqui]
		% Exemplo: 1 & 10 & texto errado & texto correto\\
		15 & 8 & conceito preliminar & conceito revisado \\
		\hline
	\end{tabular}
\end{table}
